\section{Results}
\label{sec:results}

All numbers are generated from run \texttt{\runid} by
\texttt{scripts/make\_paper\_assets.py} and \texttt{scripts/headline\_results.py}; none
is transcribed by hand. Verdicts on every pre-registered hypothesis, including those not
supported, appear in Table~\ref{tab:hypotheses}. Post-hoc analyses are labelled as such
and are itemised in the deviations section of the released analysis plan.

\subsection{Benchmark composition}

Table~\ref{tab:datasets} summarises what was evaluated: \nseries{} series and
\ntestinstances{} test instances across synthetic families with ground-truth structure
and real datasets without it.

\begin{table}[htbp]
\centering
\caption{Benchmark composition. Ground-truth structural labels exist only for synthetic
families; for real datasets the failure modes are derived rather than known.}
\label{tab:datasets}
\small
\fittable{\begin{tabular}{lrrrrr}
\toprule
Source & Family / dataset & Series & Instances & Median length & Ground-truth labels \\
\midrule
real & real & 158 & 2671 & 234 & no \\
synthetic & anomaly contaminated & 30 & 540 & 132 & yes \\
synthetic & ar stationary & 30 & 540 & 132 & yes \\
synthetic & gradual drift & 30 & 540 & 132 & yes \\
synthetic & level shift & 30 & 540 & 132 & yes \\
synthetic & missing blocks & 30 & 540 & 132 & yes \\
synthetic & multivariate var & 30 & 540 & 132 & yes \\
synthetic & near unit root & 30 & 540 & 132 & yes \\
synthetic & regime switching & 30 & 540 & 132 & yes \\
synthetic & seasonal multiple & 30 & 540 & 132 & yes \\
synthetic & seasonal single & 30 & 540 & 132 & yes \\
synthetic & short history & 30 & 60 & 35 & yes \\
synthetic & trend damped & 30 & 540 & 132 & yes \\
synthetic & trend linear & 30 & 540 & 132 & yes \\
synthetic & variance shift & 30 & 540 & 132 & yes \\
\bottomrule
\end{tabular}
}
\end{table}

\subsection{A dissociation between point and probabilistic accuracy}
\label{sec:headline}

The central result is that the answer to ``does self-correction help?'' depends entirely
on which accuracy is being measured, and that the two answers point in opposite
directions on identical instances with identical actions.

\paragraph{The unit of inference.}
Every effect below is a paired difference between per-series means, with a percentile
interval from a bootstrap that resamples \emph{series}. The pre-registered plan fixes
this unit, and it matters: our \ntestinstances{} test instances come from only
\nseries{} series, and instances within a series share a generating process, a base
model, a scaling denominator and most of their history. An instance-level interval
treats them as independent and comes out several times too narrow. We report the
signed-rank $p$-value as well, but the claims rest on the intervals. Deviation~D7 in the
analysis plan records that our own first draft got this wrong, and what changed when we
fixed it.

\begin{table}[htbp]
\centering
\caption{Correction policies under point and probabilistic accuracy. Both panels use the
same instances, the same action space and the same policies; only the scored quantity
differs. Negative percentages favour the policy over never correcting. Effects are
series-level with 95\% bootstrap intervals over series; $\dagger$ marks an interval
excluding zero. The win rate is the fraction of series on which the policy beats never
correcting, reported alongside the mean because the two disagree here.}
\label{tab:headline}
\small
\fittable{\begin{tabular}{llrrrrrr}
\toprule
Metric & Policy & Acts & Series mean & $\Delta$ & 95\% CI & Win rate & $p$ \\
\midrule
Point (MASE) & never correct & 0.00 & 0.9957 & 0.00\% & [0.00, 0.00] & 50.0\% & --- \\
Point (MASE) & rgc calibrated & 0.02 & 0.9920 & -0.38\% & [-0.97, 0.04] & 50.2\% & 0.3703 \\
Point (MASE) & always ensemble & 1.00 & 0.9491 & -4.69\% & [-11.75, 0.60] & 60.0\% & $4.2 \times 10^{-8}$ \\
Point (MASE) & always correct & 1.00 & 1.0766 & 8.13\%$^\dagger$ & [2.73, 15.99] & 45.8\% & 0.0003 \\
Point (MASE) & best fixed & 0.90 & 1.0770 & 8.16\%$^\dagger$ & [2.76, 16.05] & 45.8\% & 0.0003 \\
Point (MASE) & random correct & 0.30 & 1.0368 & 4.13\%$^\dagger$ & [2.12, 6.46] & 42.6\% & $7.0 \times 10^{-8}$ \\
Probabilistic (CRPS) & never correct & 0.00 & 0.8466 & 0.00\% & [0.00, 0.00] & 50.0\% & --- \\
Probabilistic (CRPS) & rgc calibrated & 0.97 & 0.8123 & -4.04\% & [-14.47, 2.87] & 61.0\% & $5.4 \times 10^{-7}$ \\
Probabilistic (CRPS) & always ensemble & 1.00 & 0.7359 & -13.07\%$^\dagger$ & [-22.85, -7.27] & 65.4\% & $1.2 \times 10^{-21}$ \\
Probabilistic (CRPS) & always correct & 1.00 & 0.8915 & 5.31\% & [-0.38, 12.96] & 54.8\% & 0.0406 \\
Probabilistic (CRPS) & best fixed & 0.94 & 0.8925 & 5.43\% & [-0.25, 13.05] & 53.5\% & 0.0719 \\
Probabilistic (CRPS) & random correct & 0.30 & 0.8825 & 4.25\%$^\dagger$ & [2.24, 6.58] & 42.1\% & $6.9 \times 10^{-8}$ \\
\bottomrule
\end{tabular}
}
\end{table}

\paragraph{Point accuracy: nothing reliably helps, and indiscriminate correction
reliably hurts.}
No policy establishes an improvement over accepting the initial forecast. The
unconditional static ensemble comes closest at \maseAlwaysEnsembleVsNeverPct\%, but its
interval [\maseAlwaysEnsembleVsNeverCIlo, \maseAlwaysEnsembleVsNeverCIhi]\% includes
zero; our calibrated adaptive policy is statistically indistinguishable from doing
nothing, which it achieves largely by almost never acting. The clear effects on this
metric are in the other direction: correcting every instance costs
\maseAlwaysCorrectVsNeverPct\% ([\maseAlwaysCorrectVsNeverCIlo,
\maseAlwaysCorrectVsNeverCIhi]\%) and a rate-matched random corrector
\maseRandomCorrectVsNeverPct\% ([\maseRandomCorrectVsNeverCIlo,
\maseRandomCorrectVsNeverCIhi]\%), both intervals excluding zero. Correction is not a
free option on point accuracy: applied without a reliable trigger it is a measurable
liability.

\paragraph{Probabilistic accuracy: one policy works, and it is the static one.}
The unconditional ensemble reduces series-mean CRPS by
\crpsAlwaysEnsembleVsNeverAbs\% ([\crpsAlwaysEnsembleVsNeverAbsCIlo,
\crpsAlwaysEnsembleVsNeverAbsCIhi]\%), winning on
\crpsAlwaysEnsembleVsNeverWinRate\% of series. It is the only arm on either metric whose
interval excludes zero in the helpful direction. Our calibrated adaptive gate moves in
the same direction (\crpsRgcCalibratedVsNeverPct\%) and wins on
\crpsRgcCalibratedVsNeverWinRate\% of series, but its interval
[\crpsRgcCalibratedVsNeverCIlo, \crpsRgcCalibratedVsNeverCIhi]\% includes zero, so we do
not claim it improves on never correcting. CRPS is a strictly proper scoring rule, so
this improvement cannot be manufactured by widening intervals: a point we verify
directly in Section~\ref{sec:overconfidence}, where the arm that \emph{does} widen
intervals achieves better coverage and substantially worse CRPS.

\paragraph{Was the ensemble picked on test?}
\label{sec:staticcheck}
The static ensemble was not among the pre-registered arms, so declaring it the winner
after seeing ten static actions on test would be the selection problem again. We
checked by choosing the single static action with the lowest mean error on
validation origins, blind to test, and applying it everywhere. On CRPS, validation
selects the ensemble, and its test effect is the one reported above. The CRPS result is
therefore realizable, not a post-hoc pick. On MASE, validation selects
\texttt{\spMaseValChoice}, which is significantly \emph{harmful} on test
(\spMaseValChoiceEffPct\%, [\spMaseValChoiceEffCIlo, \spMaseValChoiceEffCIhi]\%).
So on point accuracy even a validated choice of static action does not help, which is
consistent with the ensemble's own MASE interval including zero.

Because 420 of the 578 series come from our own generators, we also repeated the
comparison on each source alone. The CRPS improvement holds on the \spCrpsEnsRealN{} real
series by themselves (\spCrpsEnsRealPct\%, [\spCrpsEnsRealCIlo, \spCrpsEnsRealCIhi]\%)
and on the synthetic ones (\spCrpsEnsSynPct\%, [\spCrpsEnsSynCIlo,
\spCrpsEnsSynCIhi]\%), so it does not depend on generator design.

\paragraph{Rank and magnitude disagree, and we report both.}
For two comparisons the signed-rank test is emphatic while the bootstrap interval is
not: the ensemble beats never-correcting on \maseAlwaysEnsembleVsNeverWinRate\% of series
on MASE, and the calibrated gate on \crpsRgcCalibratedVsNeverWinRate\% of series on CRPS,
yet neither interval excludes zero. Both facts are real. A policy can win more often than
it loses while the sizes of its wins and losses fail to establish a mean improvement, which is the expected situation for a heavy-tailed loss, and the reason
Table~\ref{tab:headline} carries the win rate next to the effect rather than either
alone.

\subsection{Why: the base forecasts are over-confident}
\label{sec:overconfidence}

The mechanism is directly measurable. Nominal 80\% prediction intervals from the base
forecasts achieve substantially less than 80\% empirical coverage, and nominal 95\%
intervals likewise. Corrective actions repair the predictive distribution while leaving
its centre almost unchanged, which is exactly the pattern that a point-error metric
cannot detect and a probabilistic one can.

\paragraph{The probabilistic gain is not coverage inflation.}
An obvious objection to any claim of the form ``correction improves probabilistic
accuracy'' is that intervals can always be widened until coverage looks good, so the
result might measure nothing but timidity. Two things rule that out here. First, our
headline probabilistic metric is CRPS, a strictly proper scoring rule: widening a
predictive distribution beyond its correct dispersion increases CRPS. Coverage is
reported as mechanism, never as the outcome. Second, the benchmark contains a direct
empirical test of exactly this, because one of the ten actions does nothing else. The
\texttt{abstain} action falls back to a seasonal-naive forecast with widened intervals;
it attains the \emph{best} empirical coverage of any action in the space
(Table~\ref{tab:coverage}), essentially hitting the nominal 80\% level, with intervals
roughly $2.4\times$ wider than the base forecast's. Its CRPS is $51\%$ \emph{worse} than
accepting, and its MASE $43\%$ worse. An action that games coverage is therefore
severely penalised by the metric we actually report, which is what a proper scoring rule
is for. The ensemble improves coverage and CRPS together; abstention improves coverage
and destroys CRPS. Only the first is evidence of a better predictive distribution.

\begin{table}[htbp]
\centering
\caption{Interval coverage and CRPS by action, on the instances where each action
applies. CRPS is paired against accepting on those same instances. The base forecast is
badly over-confident; the ensemble repairs coverage and CRPS together. \texttt{abstain}
is the control for the objection that probabilistic gains can be manufactured by
widening intervals: it achieves the best coverage in the action space with intervals
roughly $2.4\times$ wider, and is penalised heavily on CRPS for it.}
\label{tab:coverage}
\small
\fittable{\begin{tabular}{lrrrrrr}
\toprule
Action & Cov.\ @80\% & Cov.\ @95\% & Calib.\ err. & Width @80\% & CRPS & CRPS $\Delta$ \\
\midrule
accept & 0.650 & 0.787 & 0.313 & 238.2 & 0.8621 & 0.0\% \\
ensemble & 0.766 & 0.906 & 0.078 & 292.6 & 0.7510 & -13.0\% \\
recalibrate intervals & 0.693 & 0.856 & 0.201 & 297.1 & 0.8453 & -2.9\% \\
abstain & 0.796 & 0.898 & 0.056 & 571.0 & 1.3038 & 51.2\% \\
damp trend & 0.572 & 0.718 & 0.460 & 238.2 & 1.0291 & 19.4\% \\
\bottomrule
\end{tabular}
}
\end{table}

\begin{figure}[htbp]
\centering
\includegraphics[width=0.8\textwidth]{fig_calibration}
\caption{Calibration of the failure-risk estimator, with expected calibration error.}
\label{fig:calibration}
\end{figure}

\subsection{What the corrections did}

Two rows carry the argument. \texttt{always\_correct} intervenes on every
instance and helps on \helpedAlwaysCorrect\% of them while harming
\harmedAlwaysCorrect\%, which is how a policy can be busy and still lose: the
harmful interventions are individually larger than the helpful ones. Our calibrated
gate intervenes on \rateRgc\% of instances and inverts that ratio, helping on
\helpedRgc\% of the instances it touches against \harmedRgc\% harmed, and its
oracle regret (\regretRgc{}) is the lowest of any acting policy.

That is the shape of the result throughout this paper. Restraint is what the gate
buys, not accuracy: it selects better than chance among the instances it acts on, and
still does not beat declining to act at all, because the actions available to it give
back more on the instances it gets wrong than they win on the instances it gets right.
Section~\ref{sec:decomposition} takes that apart.


Aggregate accuracy conceals the mechanism. Table~\ref{tab:decision} reports
decision-level quality, which only a counterfactual benchmark can produce: how often each
policy intervened, how often intervention helped rather than harmed, and the regret of
its choices against the oracle.

\begin{table}[htbp]
\centering
\caption{Decision quality. ``Helped'' and ``harmed'' are computed among corrected
instances only, so a policy can intervene frequently and still lose to never correcting.}
\label{tab:decision}
\small
\fittable{\begin{tabular}{lrrrrrrr}
\toprule
Policy & Corr. rate & Precision & Recall & F1 & Helped & Harmed & Regret \\
\midrule
never correct & 0.00 & --- & 0.000 & --- & --- & --- & 0.289 \\
rgc & 0.20 & 0.825 & 0.279 & 0.417 & 0.510 & 0.399 & 0.294 \\
random correct & 0.21 & 0.570 & 0.202 & 0.298 & 0.299 & 0.436 & 0.317 \\
threshold single & 0.29 & 0.606 & 0.296 & 0.398 & 0.188 & 0.268 & 0.350 \\
always correct & 1.00 & 0.592 & 1.000 & 0.744 & 0.264 & 0.292 & 0.378 \\
best fixed & 0.90 & 0.606 & 0.923 & 0.731 & 0.293 & 0.323 & 0.378 \\
\bottomrule
\end{tabular}
}
\end{table}

\begin{figure}[htbp]
\centering
\includegraphics[width=0.8\textwidth]{fig_correction_outcomes}
\caption{Outcome of every intervention, by policy. What matters is not how often a policy
acts but how often acting was warranted.}
\label{fig:outcomes}
\end{figure}

\subsection{Three measurement pathologies}
\label{sec:pathologies}

Each of the following concealed the effect above, and we report them because any
evaluation of correction will encounter them.

\paragraph{Ratio gains explode.}
Expressing an action's value as $(\text{accept} - \text{action}) / \text{accept}$ divides
by an error that is sometimes near zero. On that scale the correlation between an
action's observed gain on one window and its realised gain on the next is $0.006$, with
residual standard deviations above $1000$: no usable signal. On the log scale
$\log(\text{accept}/\text{action})$, the identical data gives a correlation of $0.43$
with residual standard deviations below $0.7$, and a clear per-action structure. The
conclusion changes from ``correction is unpredictable'' to ``correction is predictable
for some actions and not others''.

\paragraph{The arithmetic mean describes the tail.}
MASE on this benchmark has a median near $0.6$ and a maximum near $190$; the worst 1\% of
instances carry roughly a fifth of the total error mass. Comparisons on the mean are
therefore comparisons of a few dozen instances. We report mean, offset geometric mean and
median throughout. Absolute geometric means depend on the additive offset used to handle
exact zeros; relative comparisons do not, and only relative comparisons are interpreted.

\paragraph{Oracle headroom is mostly selection bias.}
The unrestricted oracle, the minimum over all actions on each instance, advertises
large available improvement. But it selects the best of roughly ten noisy outcomes on a
single realisation, so much of that margin is the winner's curse. A \emph{realizable}
oracle, choosing the best action on validation origins and applying it to test, captures
almost none of it (Figure~\ref{fig:oracle}). Quoting the unrestricted figure would
overstate what any policy, including a perfect one, could attain.

\begin{figure}[htbp]
\centering
\includegraphics[width=0.8\textwidth]{fig_oracle_gap}
\caption{Accept, competing policies, and the two oracles. The unrestricted oracle is an
unattainable upper bound inflated by selection; the realizable oracle is attainable.}
\label{fig:oracle}
\end{figure}

\subsection{How much headroom is there? Most of the oracle gap is selection}
\label{sec:decomposition}

A common pattern in agentic forecasting systems is to diagnose a failure mode and apply
a repair matched to it \citep{tsscientist2025,nexus2026,castflow2026}. On synthetic
instances the generator's failure mode is known, so we can grant that pattern a perfect
diagnosis and ask what it is worth. Section~\ref{sec:headline} showed that the policies
we built do not beat a static ensemble; here we ask how much there was to gain at all.

\paragraph{Scope.}
This tests our ten-action space and our hand-designed mapping from failure mode to
candidate repairs (Section~\ref{sec:actionspace}). It is not a test of any published
system, none of which we re-implemented, and a richer action space or a learned mapping
could give a different answer.

\paragraph{Four stages, and two ways to select.}
We compare accepting the forecast; a \emph{canonical fix} that applies the first
repair our mapping assigns to the true failure mode; the best of that mode's two or
three matched candidates; and the best of all ten actions. The last two involve a
choice, and the choice can be made in two ways. Selecting on the \emph{test} outcomes
gives a ceiling that no policy could reach. Selecting on \emph{validation} origins and
scoring on held-out test origins, with the rule the realizable oracle uses (per series
and horizon, the action with the lowest mean validation error), gives what a policy
with perfect per-series knowledge could actually achieve. Table~\ref{tab:regretdecomp}
reports both, with series-level bootstrap intervals throughout.

\begin{table}[htbp]
\centering
\caption{The decomposition selected two ways, on \rdMaseNseries{} synthetic series with
ground-truth diagnosis. Rows 1--3 are percentage changes relative to accepting; rows 4--5
are differences between rows, in percentage points. Intervals are 95\% series-level
bootstrap; $\dagger$ marks one excluding zero. The canonical fix involves no
selection, so its two columns are identical. The left column minimises over test
outcomes, and the gap between the columns is the winner's curse.}
\label{tab:regretdecomp}
\small
\fittable{\begin{tabular}{llrr}
\toprule
Metric & Quantity & Selected on test (ceiling) & Selected on validation \\
\midrule
MASE & Canonical fix vs accept & 2.43~[-0.50, 5.33] & 2.43~[-0.50, 5.33] \\
MASE & Best of matched candidates vs accept & -15.12$^\dagger$~[-17.60, -12.80] & 5.33~[-3.20, 19.19] \\
MASE & Best of all actions vs accept (headroom) & -31.24$^\dagger$~[-33.46, -29.10] & 4.45$^\dagger$~[0.16, 9.95] \\
MASE & First-choice loss (pp) & 17.55$^\dagger$~[14.79, 20.50] & -2.90~[-16.62, 5.49] \\
MASE & Restriction loss (pp) & 16.12$^\dagger$~[14.49, 17.87] & 0.88~[-5.13, 9.89] \\
CRPS & Canonical fix vs accept & 3.59$^\dagger$~[0.03, 7.26] & 3.59$^\dagger$~[0.03, 7.26] \\
CRPS & Best of matched candidates vs accept & -19.52$^\dagger$~[-22.05, -17.01] & 3.76~[-6.54, 20.33] \\
CRPS & Best of all actions vs accept (headroom) & -33.42$^\dagger$~[-35.77, -31.11] & -0.80~[-4.24, 2.79] \\
CRPS & First-choice loss (pp) & 23.11$^\dagger$~[19.67, 26.83] & -0.17~[-16.69, 10.45] \\
CRPS & Restriction loss (pp) & 13.90$^\dagger$~[12.28, 15.70] & 4.56~[-4.44, 19.25] \\
\bottomrule
\end{tabular}
}
\end{table}

\paragraph{Selected on test, the headroom looks large; selected on validation, it is
gone.}
Minimising over test outcomes, the best action per instance improves MASE by
\rdMaseCeilTotalAbs\% over accepting, and the gap decomposes into two sizeable losses
(\rdMaseCeilFirst{} and \rdMaseCeilRestrict{} percentage points). That is the
decomposition an earlier version of this paper reported. It is an instance of the
selection-inflated oracle described in Section~\ref{sec:pathologies}: the expected
minimum of ten noisy outcomes is lower than the minimum of three, and the minimum of
three is lower than a single draw, even when no action carries any real signal.

Made on validation, the same choices recover nothing. The best action per series
changes MASE by \rdMaseRealTotal\% ([\rdMaseRealTotalCIlo, \rdMaseRealTotalCIhi]\%),
worse than accepting, and CRPS by \rdCrpsRealTotal\% ([\rdCrpsRealTotalCIlo,
\rdCrpsRealTotalCIhi]\%). Neither named loss is distinguishable from zero. Between
\rdCrpsCurseShare\% and \rdMaseCurseShare\% of the apparent headroom was selection.
Per-series repair choices do not transfer from earlier origins to later ones on this
benchmark, so there is no realizable headroom for a per-series policy to lose.

A pre-registered arm reaches the same answer independently. \texttt{best\_fixed},
fixed in the analysis plan before any result was seen, is exactly this realizable oracle
applied to all \nseries{} series, synthetic and real. It is significantly worse than
never correcting on MASE (\maseBestFixedVsNeverPct\%, [\maseBestFixedVsNeverCIlo,
\maseBestFixedVsNeverCIhi]\%). So the result does not rest on an analysis written after a
reviewer's objection; it was already in the pre-registered results, unremarked.

\paragraph{What survives.}
Three things, none of which depends on a test-selected quantity. First, under a perfect
diagnosis the canonical repair does not help: \rdMaseCanonical\% on MASE
([\rdMaseCanonicalCIlo, \rdMaseCanonicalCIhi]\%) and \rdCrpsCanonical\% on CRPS
([\rdCrpsCanonicalCIlo, \rdCrpsCanonicalCIhi]\%), the latter a significant
degradation. Second, choosing the best repair per series on validation does not help
either. Third, a single static ensemble applied everywhere does help on CRPS
(Section~\ref{sec:headline}), and it is the action validation selects when asked for
one static policy (Section~\ref{sec:staticcheck}). Choosing per series loses to
choosing once. That is the forecast-combination puzzle
\citep{stockwatson2004,smithwallis2009,claeskens2016} reappearing at the level of repair
selection: estimating which repair suits each series adds more variance than the
specificity recovers.

This also says something about how correction policies are usually evaluated. Oracle
regret against a test-selected oracle is the standard yardstick, and on this benchmark
almost all of what it measures is selection. A policy that closes half of such a gap
may have closed none of the realizable one.

% The LLM-arm results render only when scripts/llm_arms.py has actually produced
% them. Without that file the manuscript states plainly that the arms were not run,
% instead of rendering prose about results that do not exist. See ANALYSIS_PLAN.md D8.
\IfFileExists{tables/table_llm_arms.tex}{\subsection{Letting a language model decide}
\label{sec:llmarms}

Everything so far concerns policies whose decisions are made by arithmetic. But the
systems this paper is about delegate the decision to a language model, and a critique of
that architecture that never runs it would not be worth much. This section runs it.

Three arms share the counterfactual tensor with every policy above, so the outcome of
each action is identical and the only thing that varies is who chooses.
\texttt{llm\_single} sees the diagnostic evidence once and picks an action, the
Self-Refine pattern. \texttt{llm\_critique\_gate} sees the same evidence but first judges
\emph{whether} the forecast needs revising at all: the reflection step whose value is
this paper's question, and the mechanism \citet{huang2024selfcorrect} and
\citet{kamoi2024selfcorrect} found unreliable in reasoning tasks.
\texttt{rgc\_llm\_router} hands the model only the \emph{which}, keeping our quantitative
gate for the \emph{whether}, which separates the value of gating from the value of
routing. All three receive exactly the evidence our own gate receives: the comparison is
about how the decision is made, not about who saw more.

\paragraph{How this section relates to the previous one.}
Section~\ref{sec:decomposition} found no realizable per-series headroom: choosing
repairs on validation does not beat accepting. That sets the bar for a language model
in the decision loop. It cannot win by choosing repairs cleverly, because there is
little to win; it can only avoid losing by declining to act. The question this section
asks is therefore whether a model restrains itself, and the answer is that it does not.

\paragraph{Three model families, because one would not settle it.}
A single model failing to beat never-correcting invites the obvious reply that we picked
a weak one. The arms therefore run on unrelated families: \texttt{gemini-3.5-flash-lite},
a small hosted proprietary model, on one instance from each of the \nseries{} series;
\texttt{qwen3.8-27b}, an open-weight model of different lineage served by a different
provider, on 140 series; and \texttt{gpt-oss-20b}, an open-weight \emph{reasoning} model
run at medium effort, on \reasonSeries{} series. Free-tier quotas set the coverage and they meter differently ---
one provider caps requests per day, the other tokens per day --- so the series count is
given per row in Table~\ref{tab:llmarms}, and the smaller family covers the two arms that
carry the argument rather than all three. Every non-LLM arm is re-scored within its own
family's subsample, so rows are comparable down a block but not across blocks.
Temperature is 0 throughout and every response is cached, so the experiment replays
without an API key.

The first two are not reasoning models and receive no prompt engineering: the evidence
block is a fixed template and each decision is a single zero-shot turn. The third
deliberates explicitly, and a fourth condition described below shows a model worked
examples drawn from the validation split. What the families establish together is
therefore not a claim about language models in general, but that the result is not an
artefact of one model's idiosyncrasies, of an absence of deliberation, or of the
zero-shot setting, which are the three objections a single-family null cannot answer.

\begin{table}[htbp]
\centering
\caption{LLM correction arms against the same counterfactual outcomes, for two model
families. Non-LLM arms are re-scored within each family's own subsample, so rows are
comparable within a family block but not across blocks; the covered series count is given
per row. Effects are series-level differences from never correcting with 95\% bootstrap
intervals; $\dagger$ marks an interval excluding zero. Negative favours the policy.
``LLM calls'' is the live provider call count, since a correction stage that does not
help is still not free.}
\label{tab:llmarms}
\small
\fittable{\begin{tabular}{lllrrrrrr}
\toprule
Metric & Model & Policy & Series & Series mean & $\Delta$ & 95\% CI & Win rate & LLM calls \\
\midrule
Point (MASE) & Gemini 3.5 Flash-Lite & never correct & 578 & 1.2220 & 0.00\% & [0.00, 0.00] & 50.0\% & 0 \\
Point (MASE) & Gemini 3.5 Flash-Lite & always ensemble & 578 & 1.1520 & -5.73\%$^\dagger$ & [-9.54, -1.99] & 53.9\% & 0 \\
Point (MASE) & Gemini 3.5 Flash-Lite & rgc calibrated & 578 & 1.2196 & -0.19\% & [-0.72, 0.26] & 50.3\% & 0 \\
Point (MASE) & Gemini 3.5 Flash-Lite & llm single & 578 & 1.3243 & 8.37\%$^\dagger$ & [3.29, 13.94] & 42.2\% & 0 \\
Point (MASE) & Gemini 3.5 Flash-Lite & llm critique gate & 578 & 1.3136 & 7.50\%$^\dagger$ & [2.73, 12.89] & 42.8\% & 288 \\
Point (MASE) & Gemini 3.5 Flash-Lite & rgc llm router & 578 & 1.2814 & 4.87\%$^\dagger$ & [0.95, 10.06] & 47.5\% & 258 \\
Point (MASE) & Gemini 3.5 Flash-Lite & always correct & 578 & 1.2920 & 5.73\% & [-0.48, 14.43] & 47.1\% & 0 \\
Probabilistic (CRPS) & Gemini 3.5 Flash-Lite & never correct & 578 & 1.0629 & 0.00\% & [0.00, 0.00] & 50.0\% & 0 \\
Probabilistic (CRPS) & Gemini 3.5 Flash-Lite & always ensemble & 578 & 0.9503 & -10.59\%$^\dagger$ & [-15.14, -6.70] & 53.8\% & 0 \\
Probabilistic (CRPS) & Gemini 3.5 Flash-Lite & rgc calibrated & 578 & 1.0064 & -5.31\%$^\dagger$ & [-9.36, -1.91] & 53.5\% & 0 \\
Probabilistic (CRPS) & Gemini 3.5 Flash-Lite & llm single & 578 & 1.1530 & 8.48\%$^\dagger$ & [3.00, 14.55] & 43.9\% & 0 \\
Probabilistic (CRPS) & Gemini 3.5 Flash-Lite & llm critique gate & 578 & 1.1493 & 8.13\%$^\dagger$ & [3.17, 13.92] & 44.3\% & 288 \\
Probabilistic (CRPS) & Gemini 3.5 Flash-Lite & rgc llm router & 578 & 1.1015 & 3.63\%$^\dagger$ & [0.48, 8.04] & 48.7\% & 258 \\
Probabilistic (CRPS) & Gemini 3.5 Flash-Lite & always correct & 578 & 1.0979 & 3.29\% & [-4.92, 13.64] & 51.6\% & 0 \\
Point (MASE) & Qwen3.8-27B (Groq) & never correct & 140 & 0.8936 & 0.00\% & [0.00, 0.00] & 50.0\% & 0 \\
Point (MASE) & Qwen3.8-27B (Groq) & always ensemble & 140 & 0.8892 & -0.50\% & [-5.17, 4.69] & 51.4\% & 0 \\
Point (MASE) & Qwen3.8-27B (Groq) & rgc calibrated & 140 & 0.8850 & -0.96\% & [-2.85, 0.00] & 50.7\% & 0 \\
Point (MASE) & Qwen3.8-27B (Groq) & llm single & 140 & 0.9400 & 5.19\%$^\dagger$ & [0.86, 10.19] & 46.8\% & 0 \\
Point (MASE) & Qwen3.8-27B (Groq) & llm critique gate & 140 & 0.9458 & 5.85\%$^\dagger$ & [1.45, 10.84] & 45.7\% & 11 \\
Probabilistic (CRPS) & Qwen3.8-27B (Groq) & never correct & 140 & 0.7513 & 0.00\% & [0.00, 0.00] & 50.0\% & 0 \\
Probabilistic (CRPS) & Qwen3.8-27B (Groq) & always ensemble & 140 & 0.7168 & -4.59\% & [-10.64, 0.96] & 52.1\% & 0 \\
Probabilistic (CRPS) & Qwen3.8-27B (Groq) & rgc calibrated & 140 & 0.7097 & -5.54\%$^\dagger$ & [-10.73, -0.69] & 57.1\% & 0 \\
Probabilistic (CRPS) & Qwen3.8-27B (Groq) & llm single & 140 & 0.7941 & 5.70\% & [-0.48, 12.33] & 48.6\% & 0 \\
Probabilistic (CRPS) & Qwen3.8-27B (Groq) & llm critique gate & 140 & 0.7892 & 5.04\% & [-1.22, 11.81] & 50.7\% & 11 \\
Point (MASE) & GPT-OSS-20B reasoning (Groq) & never correct & 60 & 1.0447 & 0.00\% & [0.00, 0.00] & 50.0\% & 0 \\
Point (MASE) & GPT-OSS-20B reasoning (Groq) & always ensemble & 60 & 1.0974 & 5.04\% & [-9.06, 24.71] & 48.3\% & 0 \\
Point (MASE) & GPT-OSS-20B reasoning (Groq) & rgc calibrated & 60 & 1.0454 & 0.06\% & [-0.16, 0.34] & 50.0\% & 0 \\
Point (MASE) & GPT-OSS-20B reasoning (Groq) & llm single & 60 & 1.1024 & 5.52\% & [-2.92, 14.61] & 50.0\% & 60 \\
Point (MASE) & GPT-OSS-20B reasoning (Groq) & llm critique gate & 60 & 1.1034 & 5.62\% & [-3.77, 15.32] & 47.5\% & 60 \\
Probabilistic (CRPS) & GPT-OSS-20B reasoning (Groq) & never correct & 60 & 0.9179 & 0.00\% & [0.00, 0.00] & 50.0\% & 0 \\
Probabilistic (CRPS) & GPT-OSS-20B reasoning (Groq) & always ensemble & 60 & 0.8967 & -2.31\% & [-16.91, 13.54] & 53.3\% & 0 \\
Probabilistic (CRPS) & GPT-OSS-20B reasoning (Groq) & rgc calibrated & 60 & 0.9341 & 1.76\% & [-3.18, 7.54] & 51.7\% & 0 \\
Probabilistic (CRPS) & GPT-OSS-20B reasoning (Groq) & llm single & 60 & 0.9776 & 6.49\% & [-2.30, 15.67] & 42.5\% & 60 \\
Probabilistic (CRPS) & GPT-OSS-20B reasoning (Groq) & llm critique gate & 60 & 0.9805 & 6.81\% & [-1.72, 15.52] & 43.3\% & 60 \\
Point (MASE) & Gemini 3.5 Flash-Lite, 8-shot & never correct & 60 & 1.0447 & 0.00\% & [0.00, 0.00] & 50.0\% & 0 \\
Point (MASE) & Gemini 3.5 Flash-Lite, 8-shot & always ensemble & 60 & 1.0974 & 5.04\% & [-9.06, 24.71] & 48.3\% & 0 \\
Point (MASE) & Gemini 3.5 Flash-Lite, 8-shot & rgc calibrated & 60 & 1.0454 & 0.06\% & [-0.16, 0.34] & 50.0\% & 0 \\
Point (MASE) & Gemini 3.5 Flash-Lite, 8-shot & llm single & 60 & 1.1142 & 6.65\% & [-1.25, 15.77] & 45.8\% & 0 \\
Point (MASE) & Gemini 3.5 Flash-Lite, 8-shot & llm critique gate & 60 & 1.3462 & 28.85\% & [-0.49, 80.30] & 46.7\% & 0 \\
Probabilistic (CRPS) & Gemini 3.5 Flash-Lite, 8-shot & never correct & 60 & 0.9179 & 0.00\% & [0.00, 0.00] & 50.0\% & 0 \\
Probabilistic (CRPS) & Gemini 3.5 Flash-Lite, 8-shot & always ensemble & 60 & 0.8967 & -2.31\% & [-16.91, 13.54] & 53.3\% & 0 \\
Probabilistic (CRPS) & Gemini 3.5 Flash-Lite, 8-shot & rgc calibrated & 60 & 0.9341 & 1.76\% & [-3.18, 7.54] & 51.7\% & 0 \\
Probabilistic (CRPS) & Gemini 3.5 Flash-Lite, 8-shot & llm single & 60 & 0.9944 & 8.33\%$^\dagger$ & [0.49, 16.66] & 40.0\% & 0 \\
Probabilistic (CRPS) & Gemini 3.5 Flash-Lite, 8-shot & llm critique gate & 60 & 1.1299 & 23.09\%$^\dagger$ & [1.96, 58.52] & 39.2\% & 0 \\
Point (MASE) & Gemini 3.5 Flash-Lite, 0-shot (matched) & never correct & 60 & 1.0447 & 0.00\% & [0.00, 0.00] & 50.0\% & 0 \\
Point (MASE) & Gemini 3.5 Flash-Lite, 0-shot (matched) & always ensemble & 60 & 1.0974 & 5.04\% & [-9.06, 24.71] & 48.3\% & 0 \\
Point (MASE) & Gemini 3.5 Flash-Lite, 0-shot (matched) & rgc calibrated & 60 & 1.0454 & 0.06\% & [-0.16, 0.34] & 50.0\% & 0 \\
Point (MASE) & Gemini 3.5 Flash-Lite, 0-shot (matched) & llm single & 60 & 1.1330 & 8.45\%$^\dagger$ & [3.00, 14.72] & 39.2\% & 50 \\
Point (MASE) & Gemini 3.5 Flash-Lite, 0-shot (matched) & llm critique gate & 60 & 1.1407 & 9.19\%$^\dagger$ & [1.69, 17.60] & 38.3\% & 50 \\
Probabilistic (CRPS) & Gemini 3.5 Flash-Lite, 0-shot (matched) & never correct & 60 & 0.9179 & 0.00\% & [0.00, 0.00] & 50.0\% & 0 \\
Probabilistic (CRPS) & Gemini 3.5 Flash-Lite, 0-shot (matched) & always ensemble & 60 & 0.8967 & -2.31\% & [-16.91, 13.54] & 53.3\% & 0 \\
Probabilistic (CRPS) & Gemini 3.5 Flash-Lite, 0-shot (matched) & rgc calibrated & 60 & 0.9341 & 1.76\% & [-3.18, 7.54] & 51.7\% & 0 \\
Probabilistic (CRPS) & Gemini 3.5 Flash-Lite, 0-shot (matched) & llm single & 60 & 0.9981 & 8.73\%$^\dagger$ & [2.75, 15.62] & 38.3\% & 50 \\
Probabilistic (CRPS) & Gemini 3.5 Flash-Lite, 0-shot (matched) & llm critique gate & 60 & 1.0195 & 11.06\%$^\dagger$ & [4.34, 18.85] & 34.2\% & 50 \\
\bottomrule
\end{tabular}
}
\end{table}

\paragraph{On the primary family, every LLM arm is worse than doing nothing.}
On the primary family all
three arms degrade both metrics with intervals that exclude zero: a single pass costs
\llmMaseLlmSinglePct\% MASE ([\llmMaseLlmSingleCIlo, \llmMaseLlmSingleCIhi]\%) and
\llmCrpsLlmSinglePct\% CRPS ([\llmCrpsLlmSingleCIlo, \llmCrpsLlmSingleCIhi]\%); the
self-critique gate costs \llmMaseLlmCritiqueGatePct\% and
\llmCrpsLlmCritiqueGatePct\%, and even the router, which acts on only a seventh of
instances because our quantitative gate holds it back, costs
\llmMaseRgcLlmRouterPct\% and \llmCrpsRgcLlmRouterPct\%. Each wins on fewer than half of
series. On the same instances the static ensemble improves CRPS by
\llmCrpsAlwaysEnsembleAbs\% ([\llmCrpsAlwaysEnsembleAbsCIlo,
\llmCrpsAlwaysEnsembleAbsCIhi]\%), so the subsample is not too small to resolve a real
effect; it resolves one, in the opposite direction from the arms that reason about each
instance.

\paragraph{A discrepancy between two designs.}
In Table~\ref{tab:llmarms} the static ensemble's MASE advantage clears significance,
whereas in Table~\ref{tab:headline} the same policy's MASE interval includes zero. The
tables are computed on different data: the headline analysis uses every test origin
(\ntestinstances{} instances), while these arms use one randomly chosen origin per series,
because that is what the API budget allowed. Both are series-level and both are correct
for what they cover; the subsample simply draws fewer of the extreme origins that widen
the full interval. \textbf{The headline analysis is the primary one}, and the claim we
stand behind is the more conservative of the two: that no policy reliably improves
point accuracy. We flag the difference because quoting the subsample's narrower interval
as though it settled the point would be exactly the kind of selection this paper
criticises elsewhere.

The decision-level view explains the aggregate. When the language model chooses to
correct, it makes the forecast \emph{worse} on roughly half the instances it touches and
better on under a third: it harms about $1.8\times$ as often as it helps. The router
is the worst offender at a $64\%$ harm rate, which is initially surprising until one
notices what it means: even after a calibrated gate has correctly identified an instance
worth correcting, delegating the choice of \emph{which} repair to a language model
destroys the opportunity. That is Section~\ref{sec:decomposition}'s restriction-and-first-choice
loss appearing again, with a language model in the role of the fixed canonical mapping.
A better-informed guess is still a guess; it is not a measurement.

\paragraph{The reflection step carries almost no information.}
We expected the comparison between \texttt{llm\_single} and \texttt{llm\_critique\_gate}
to be the sharpest in this section, because they differ in exactly one thing: whether the
model is asked to judge the need for correction before choosing. It is not sharp at all.
Asked plainly whether the forecast needs revising, the model answers yes on
\rateLlmCritiqueGate\% of instances, almost exactly the rate at which the single-pass arm
chooses to act anyway, and the two arms land within a percentage point of each other on
both metrics.

Matching rates would be weak evidence on its own; two arms can intervene equally often
while selecting very different instances, one of them well. So we test discrimination
directly, against the benchmark's ground-truth \texttt{should\_correct} label
(Table~\ref{tab:discrimination}). The critique's yes/no separates warranted from
unwarranted correction at AUROC \discrimAurocCritique{}, against a chance value of 0.500.
Its selection precision is \discrimPrecisionCritique\%, where the base rate is
\discrimBaseRate\% and a coin firing at the same rate achieves \discrimPrecisionRandom\%.
And \discrimJaccard\% of the instances it elects to correct are the same instances the
uncritiqued arm corrects.

The decisive comparison is with the arm that does no critique at all, which discriminates
at AUROC \discrimAurocSingle{}. \textbf{Adding an explicit reflection step moves
discrimination by \discrimAurocDelta{}; from near-chance to near-chance.} The critique
is not merely a weak gate; it is very nearly a re-description of what the model would have
done without being asked. Both families show the same pattern
(Table~\ref{tab:discrimination}), including one where the critique step \emph{raises} the
intervention rate instead of leaving it flat, and still adds essentially nothing to
discrimination.

\paragraph{The same judgement, made arithmetically.}
The natural question is whether this decision is simply hard. It is not. Our reliability
estimator --- gradient boosting over the same diagnostic features the model was shown in
prose, fitted on validation only --- scores \discrimAurocArithmetic{} against the
\emph{same} label on the \emph{same} instances (Table~\ref{tab:discrimination}). That is
\discrimArithmeticAdvantage{} AUROC above the self-critique. We deliberately do not
express that as a multiple of the critique's distance from chance: dividing by a
near-zero baseline is the first of the measurement pathologies this paper documents
(Section~\ref{sec:pathologies}). We also
report the contrast \emph{matched}, because the triage figure quoted in
Section~\ref{sec:triage} targets a different label, and setting the two against each
other directly would not be legitimate.

\paragraph{What that contrast does and does not license.}
A binary yes/no scored as an AUROC is balanced accuracy, and nothing more. The
estimator emits a continuous score, so scoring the two the same way credits it for
ranking that a yes/no cannot express, and some of the apparent gap is an artefact of
the output type rather than of the judgement behind it. The contrast is therefore reported at matched footing in two ways, both in
Table~\ref{tab:discrimination}.

The first gives the language model the same kind of credit. Its verbalised confidence
is continuous, so it can be scored as a ranking against the same label: it reaches
AUROC \discrimAurocLlmConfidence{} against the estimator's \discrimAurocArithmetic{},
a gap of \discrimContinuousGap{}. The second removes the ranking question entirely
by forcing both to the same operating point --- the estimator thresholded to flag
exactly the fraction of instances the critique gate flags --- and compares balanced
accuracy: \discrimBalaccArithmetic{} against \discrimBalaccCritique{}, a gap of
\discrimMatchedGap{}.

The two ask different questions and both are worth keeping. (a) asks whether the
model's stated risk can \emph{rank} instances at all, independent of where a threshold
would be placed. (b) asks how much better the decision is at the rate the model
actually fires, which is the quantity a deployed gate delivers. It is worth noting that
(a) is marginally \emph{larger} than the unmatched comparison it replaces
(\discrimContinuousGap{} against \discrimArithmeticAdvantage{}), so matching was not a
route to a smaller number; only (b), which removes the ranking question entirely,
narrows the gap.

At a matched operating point the gap
is \discrimMatchedGap{}, not the \discrimArithmeticAdvantage{} a naive reading of the
first two columns would give. The artefact was real and it was worth roughly a third
of the apparent difference. What survives it is still the finding: on the same
instances, against the same label, at the same rate of intervention, a gradient-boosted
model over a few dozen scalar diagnostics decides substantially better than a language
model reasoning over the same evidence in words.

What the contrast does not license is a claim that language models cannot do this. The
estimator is fitted on the validation split and the model is not, and that asymmetry is
addressed directly below. Nor does it license a claim about relative ceilings: it
compares one fitted estimator against one prompted model, at one budget, and a
different prompt or a different model could land elsewhere. It licenses exactly one
thing, which is the thing the deployment question turns on: that the cheap
arithmetic option available to any practitioner today is the better of the two.


\begin{table}[htbp]
\centering
\caption{Does verbal self-critique carry information about whether to correct? All
columns are computed on the same instances against the same ground-truth
\texttt{should\_correct} label. ``AUROC (critique)'' scores a binary yes/no and is
therefore balanced accuracy; ``AUROC (arithmetic)'' scores a continuous estimator and
is credited for ranking, so the two are compared at matched footing in the text rather
than by subtracting these columns. $\Delta$ AUROC is the within-model quantity the
claim rests on: the same model, the same instances, deciding with and without being
asked to reflect. ``Overlap'' is the Jaccard overlap of the two intervention sets.}
\label{tab:discrimination}
\small
\fittable{\begin{tabular}{lrrrrrrrrr}
\toprule
Model & $n$ & Base rate & AUROC (critique) & AUROC (no critique) & $\Delta$ AUROC & AUROC (arithmetic) & Precision & Random & Overlap \\
\midrule
Gemini 3.5 Flash-Lite & 578 & 59.7\% & 0.532 & 0.522 & 0.010 & 0.760 & 61.9\% & 59.9\% & 82\% \\
Qwen3.8-27B (Groq) & 140 & 52.9\% & 0.568 & 0.561 & 0.008 & 0.693 & 59.0\% & 51.4\% & 75\% \\
GPT-OSS-20B reasoning (Groq) & 60 & 60.0\% & 0.493 & 0.590 & -0.097 & 0.774 & 59.6\% & 60.0\% & 80\% \\
\bottomrule
\end{tabular}
}
\end{table}

\paragraph{Why it intervenes so often: it thinks its forecasts are far worse than
they are.}
The arms are asked, in the same breath as the decision, for a number between 0 and 1
expressing how reliable they judge the current forecast to be. That is a claim about
exactly the event the reliability estimator predicts, so the two can be scored against
each other on identical instances and an identical label. This is pre-registered
hypothesis H5, and it is \textbf{supported}: the estimator's calibration error is
\calEceDiff{} lower ([\calEceDiffCIlo, \calEceDiffCIhi], bootstrapped over series),
and the same ordering holds on the Brier score.

The size and direction of the miscalibration are what matter here. Large errors occur
on \calBaseRate\% (on the \calN-instance subsample the LLM arms run on) of these instances; the model's stated risk averages
\calLlmMeanRisk{}. It believes its forecasts fail roughly two and a half times as often
as they do. That is not a detail about calibration, it is the mechanism behind the
headline: a model that thinks most of its forecasts are unreliable will conclude that
most of them need fixing, which is precisely the disposition the intervention rates
show. The reflection step is not failing to detect a problem it should have seen: it is reporting a problem that is mostly not there.

Table~\ref{tab:calibration} also contains the reason we do not rest this on calibration
error alone. A constant predictor that simply emits the base rate attains an ECE of
essentially zero while ranking nothing at all, and its Brier score (\calConstBrier{})
is nonetheless better than the language model's (\calLlmBrier{}). By the proper
scoring rule, the agent's stated confidence is worse than a constant. It is not
worthless as a \emph{ranking}: its AURC of \calLlmAurc{} beats the constant's
\calConstAurc{}, so there is some signal in which forecasts it doubts most. But that
signal is buried under a bias large enough to make the resulting decisions wrong, and a
gate reads the level, not the ranking.

That a language model's stated confidence is poorly calibrated is not itself news
\citep{kadavath2022know,lin2022verbalized,xiong2024confidence}, and we do not present it
as such. Two things here are worth separating from that literature. The direction is
inverted: the usual finding is overconfidence, a model claiming more certainty than its
accuracy warrants, whereas this model is systematically \emph{pessimistic} about the
artefact it is judging. And the consequence is operational rather than presentational.
In the settings where verbalised confidence is normally studied it is a number shown to
a reader, who may discount it; here it is wired into a control loop, where a bias of
this size does not mislead anyone so much as it manufactures roughly one unnecessary
repair for every instance that needed one.

\begin{table}[htbp]
\centering
\caption{H5: the agent's stated confidence against a calibrated estimate of the same
event. Both predict the probability of a large error, on the same \calN{} instances
and against the same label; lower is better in every column. The constant predictor is
included deliberately: it attains a perfect calibration error while ranking nothing,
which is why ECE is never read here on its own.}
\label{tab:calibration}
\small
\fittable{\begin{tabular}{lrrrr}
\toprule
Predictor & ECE & Brier & AURC & Mean $p$ \\
\midrule
LLM verbalised confidence & 0.3254 & 0.3148 & 0.1593 & 0.544 \\
Reliability estimator (calibrated) & 0.0502 & 0.1401 & 0.0883 & 0.185 \\
Constant at base rate (0.22) & 0.0000 & 0.1714 & 0.2291 & 0.220 \\
\bottomrule
\end{tabular}
}
\end{table}

\paragraph{Deliberation, and demonstration, against criteria fixed in advance.}
Two reasonable objections are still open at this point. The models run so far answer in a
single zero-shot turn, so the failure could be an absence of deliberation instead of a
property of the mechanism. And the arithmetic comparator is fitted on the validation
split while the language model is shown nothing, so the contrast above could be measuring
supervision rather than capability. Both were tested on one shared denominator ---
\reasonSeries{} series, one origin each --- with the falsification criteria written down
before the runs completed (deviation~D9). The deliberation arm is \texttt{gpt-oss-20b} at
\texttt{reasoning\_effort} \texttt{= medium}, with a completion budget large enough for
the reasoning trace and the answer; it recorded no provider failures and no unparsed
replies. The demonstration arm prepends eight worked examples drawn from validation and
labelled by the counterfactual tensor --- the same supervision, out of the same data, that
the estimator receives --- and runs against a matched zero-shot control on the identical
series, so the comparison is within-model on the same instances.

Neither closes the gap. The reasoning model degrades MASE by \reasonMaseSinglePct\%
([\reasonMaseSingleCIlo, \reasonMaseSingleCIhi]\%) in one pass and
\reasonMaseCritiquePct\% ([\reasonMaseCritiqueCIlo, \reasonMaseCritiqueCIhi]\%) with
self-critique, intervening on \reasonRateSingle\% and \reasonRateCritique\% of instances
respectively: the highest rates of any family. Eight worked examples move the
intervention rate from \zeroshotRateSingle\% to \fewshotRateSingle\% and the effect on
MASE from \zeroshotMaseSinglePct\% ([\zeroshotMaseSingleCIlo,
\zeroshotMaseSingleCIhi]\%, excluding zero) to \fewshotMaseSinglePct\%
([\fewshotMaseSingleCIlo, \fewshotMaseSingleCIhi]\%, no longer excluding zero). That is a small improvement this denominator cannot resolve, not a repair: the
two intervals overlap over almost their whole length, and the demonstration condition does
not bring the arm to parity with never correcting. The critique gate is a different story, and it runs against us. With the same eight
examples its degradation grew rather than shrank: from \zeroshotMaseCritiquePct\% to
\fewshotMaseCritiquePct\% on MASE ([\fewshotMaseCritiqueCIlo, \fewshotMaseCritiqueCIhi]\%)
and from \zeroshotCrpsCritiquePct\% to \fewshotCrpsCritiquePct\% on CRPS
([\fewshotCrpsCritiqueCIlo, \fewshotCrpsCritiqueCIhi]\%, excluding zero). The MASE
interval is wide enough that a few series drive it, but the direction is unambiguous:
showing the gate worked examples made it more harmful, not less. We do not have an
explanation we can test. What it does rule out is the reading that the critique gate
failed only for want of supervision.

The obvious question is why the demonstration condition stopped at eight examples and
\fewshotSeries{} series. The answer is budget rather than design: the prefix multiplies
prompt tokens roughly sixfold, and free-tier quotas bind on tokens for one provider and
on requests for the other. A larger demonstration budget, or in-context examples
selected per instance rather than fixed, is untested here and is the version of this
arm we would run first with more compute. What we can say is that the cheapest form of
the intervention, the one a practitioner would reach for, does not restrain the
model, and that is a narrower claim than the arm was built to make.

The pre-registered falsifiers did not fire. The primary one asked whether self-critique
would raise this model's discrimination by more than $+0.05$ AUROC over the same model
deciding without reflection. It moved by \reasonAurocDelta{}: \emph{down}, from
\reasonAurocSingle{} to \reasonAurocCritique{}, on \reasonDiscrimN{} instances where the
arithmetic estimator scored \reasonAurocArithmetic{} against the identical label. This is
the strongest form of the within-model result in the paper. Asking the one model that
actually deliberates to deliberate about its own judgement did not improve that judgement;
it made it slightly worse.

\paragraph{Given the decision to correct, is the chosen repair any good?}
The arms conflate two decisions: whether to correct, and which repair to apply. The
tensor separates them, because it holds the outcome of every action on every instance.
On the instances where an arm intervened we compare the repair it chose against a
uniform draw from the same menu on the same instance, computed exactly by averaging the
menu's outcomes. That comparison involves no selection, so it is the one the claims rest
on (Table~\ref{tab:actionchoice}).

On the primary family the choice cannot be told apart from the coin. Single-pass
selection scores \actionChoiceSinglePct\% relative to a uniform draw
([\actionChoiceSingleCIlo, \actionChoiceSingleCIhi]\%) and the critique arm
\actionChoiceCritiquePct\% ([\actionChoiceCritiqueCIlo, \actionChoiceCritiqueCIhi]\%).
Having decided to act, the model's choice of repair carries no measurable advantage over
choosing at random among the options it was shown.

The table also shows the best action in each menu. That row is a maximum over test
outcomes, and Section~\ref{sec:decomposition} shows such maxima are almost entirely
winner's curse on this benchmark, so it is context rather than a target, and no claim
uses it. The router arm's \actionChoiceRouterPct\% rests on \actionChoiceRouterN{}
instances with an interval of [\actionChoiceRouterCIlo, \actionChoiceRouterCIhi]\%; it
is underpowered and excluded from claims.

\begin{table}[htbp]
\centering
\caption{Given the decision to correct, how good is the chosen repair? MASE on the
instances where each arm intervened. The uniform draw averages the menu's outcomes
exactly and is the comparison claims rest on. The best-in-menu row is a maximum over
test outcomes, inflated by selection (Section~\ref{sec:decomposition}), and is shown
for context only. Intervals are series-level bootstrap; the router arm is underpowered.}
\label{tab:actionchoice}
\small
\fittable{\begin{tabular}{lrrr}
\toprule
 & LLM single & LLM critique gate & RGC + LLM router \\
\midrule
Instances intervened on & 361 & 367 & 80 \\
Best in menu (test-selected ceiling) & 0.6487 & 0.6273 & 0.5067 \\
\textbf{Model's chosen action} & \textbf{1.1190} & \textbf{1.0906} & \textbf{1.4109} \\
Uniform draw from same menu & 1.0806 & 1.0678 & 1.0996 \\
Worst action in menu & 2.0608 & 2.0461 & 2.5194 \\
\midrule Chosen vs.\ uniform draw & 3.55\%~[-4.07, 12.11] & 2.13\%~[-5.17, 10.66] & 28.31\%~[-0.99, 65.84] \\
\bottomrule
\end{tabular}
}
\end{table}

\paragraph{The reasoning model may choose better, but it still acts too often.}
The reasoning family points the other way on this test. Its single-pass choice scores
\reasonChoiceSinglePct\% against a uniform draw ([\reasonChoiceSingleCIlo,
\reasonChoiceSingleCIhi]\%) and its critique arm \reasonChoiceCritiquePct\%
([\reasonChoiceCritiqueCIlo, \reasonChoiceCritiqueCIhi]\%). Both intervals just
include zero at \reasonSeries{} series, so this is a direction, not an effect. If it
holds, deliberation helps with which repair to apply and not with whether to apply one;
and because this arm still intervenes on \reasonRateSingle\% of instances, its net
effect on accuracy stays negative either way.

\paragraph{The pre-registered test of H7.}
H7 asked whether a quantitative gate decides better than verbal self-critique, scored by
oracle regret. It is the one hypothesis in the plan that could not be evaluated until
these arms ran. It is \textbf{\verdictHSeven}: the calibrated gate's mean regret is
\statHSeven{} lower than the LLM critique gate's. The comparison is worth stating
carefully, because the gate achieves this largely by declining to act; its regret is
almost identical to never correcting. The finding is therefore not that quantitative
gating is good at choosing corrections. It is that it is good at \emph{not} choosing
them, and that verbal self-critique is not.

\paragraph{What this does not establish.}
Three families, a reasoning model at real effort, and a demonstration condition close the
cheapest objections but not every one. We cannot rule out that a frontier reasoning model,
multi-turn deliberation with tool access during the decision, or a prompt engineered
against this benchmark would do better; nor that a larger member of the same reasoning
family --- the one originally planned, lost to a daily token cap --- would separate from
the 20B model we ran. The reasoning and demonstration conditions cover \reasonSeries{}
series, which is enough to replicate a direction and not enough to establish an effect on
its own. Nor did we test sensitivity to prompt wording, and the pessimism result in particular
may depend on how the confidence question is phrased. What the result does establish
is narrower: the mechanism
\emph{as commonly deployed}, on the evidence commonly available to it, is not merely
unhelpful but reliably harmful; its harm resolves at the primary family's sample size, while the two smaller families show the same direction without resolving it. Self-reflection, the intervention most often proposed as the fix, did not improve the model's judgement in any family we ran.

\paragraph{The cost side.}
These are also the expensive arms. The self-critique gate alone made
\llmCallsCritiqueGate{} live provider calls, and the three arms together consumed
\llmCallsTotal{} calls and \llmPromptTokensThousands{}k prompt tokens on the primary
family. The static ensemble that outperforms all of them makes zero
provider calls and takes no per-instance decision whatsoever. Any argument for an
LLM-in-the-loop correction stage has to be made against that baseline, and on this
benchmark it cannot be made on accuracy at any price.
}{%
\subsection{Letting a language model decide}
\label{sec:llmarms}

\emph{This section is intentionally empty in the present build.} The arms that delegate
the correction decision to a language model --- a single-pass chooser, a verbal
self-critique gate, and a gate-plus-LLM-router --- are implemented, tested end to end,
and run with one command against the same counterfactual tensor as every policy above.
They are not reported here because the provider's free tier allows 500 requests per day
for the configured model and 20 for the alternative, and a run covering enough series to
support a series-level interval needs several times that. Deviation~D8 in the analysis
plan records the constraint and what was cached.

We flag this prominently instead of burying it: this manuscript argues about
architectures that delegate correction to a language model, and until these arms are
reported that argument rests on the decomposition in
Section~\ref{sec:decomposition} and on the non-LLM policies, not on direct evidence about
language-model decision-making. A reader should discount the paper's claims about verbal
self-critique accordingly.
}

\subsection{Selective prediction: a positive, deployable finding}
\label{sec:triage}

The results above are cautionary about automated correction. The failure-risk estimator
that drives the gate is nonetheless independently useful for a different and arguably
more common industrial purpose: prioritising forecasts for human review. Because the
estimator is fit and evaluated exactly as in Section~\ref{sec:decomposition} (validation
only, never test), the numbers below describe a deployable capability instead of a
ceiling.

\begin{table}[htbp]
\centering
\caption{Selective-prediction value of the failure-risk estimator (AUROC
$=\triageauroc$, base rate of large errors $=\triagebaserate\%$). Reviewing the
riskiest 10\% of forecasts catches roughly a third of all large errors at two-thirds
precision -- a $\triageliftPOneZero\times$ enrichment over reviewing at random.}
\label{tab:triage}
\small
\fittable{\begin{tabular}{lrrrr}
\toprule
Review budget & Reviewed & Precision & Recall & Lift \\
\midrule
5.0\% & 164 & 0.762 & 0.190 & 3.81$\times$ \\
10.0\% & 328 & 0.665 & 0.332 & 3.32$\times$ \\
20.0\% & 656 & 0.511 & 0.510 & 2.55$\times$ \\
30.0\% & 985 & 0.424 & 0.636 & 2.12$\times$ \\
50.0\% & 1642 & 0.322 & 0.805 & 1.61$\times$ \\
\bottomrule
\end{tabular}
}
\end{table}

This reframes what the benchmark's reliability estimation contributes: not necessarily
\emph{automated} self-correction, which our results treat with substantial scepticism,
but \emph{calibrated triage} -- flagging which forecasts warrant scrutiny, at a
quantified precision/recall operating point a practitioner can choose to match a review
budget. This is a positive, actionable, and immediately deployable use of the same
signal that the correction-policy analysis shows is hard to act on automatically.

\subsection{Failure is predictable in advance}

H3 asked whether a forecast's failure can be anticipated before the outcome is known,
using only history-side diagnostics. It can. The gradient-boosted risk estimator,
fitted on validation and evaluated on held-out test instances, reaches AUROC
\triageauroc{} against a large-error base rate of \triagebaserate\%. It is
\textbf{\verdictHThree}.

This is worth separating carefully from everything else here, because it is the one
component that works as designed and it is easy to over-read. Predicting that a
forecast will be bad is not the same as knowing what to do about it. The same
estimator drives the correction gate, and the gate does not beat never correcting
(Table~\ref{tab:headline}); it is the action that fails, not the prediction. What the
signal does support is triage, where a human decides what to do next, and that is the
subject of Section~\ref{sec:triage}.

Figure~\ref{fig:importance} ranks the signals by permutation importance: the drop
in AUROC when a single signal is shuffled. The ablations in
Section~\ref{sec:ablations} give the complementary view, measuring what each signal is
worth to the end-to-end policy rather than to the classifier, and the two do not
agree: a signal can carry information about failure and still not improve a decision
made on top of it.


\begin{figure}[htbp]
\centering
\includegraphics[width=0.62\textwidth]{fig_feature_importance}
\caption{Signals driving the correction gate, by permutation importance (drop in AUROC
when a signal is shuffled).}
\label{fig:importance}
\end{figure}

\subsection{Cost}

Accuracy is only half of a deployment decision. Table~\ref{tab:cost} prices each
policy in action evaluations per instance, the unit that dominates: every evaluation
is a re-forecast.

The ordering is unkind to adaptivity. \texttt{always\_correct} evaluates
\deltaAlwaysCorrect{} actions per instance to finish worse than evaluating none. Our
gate spends \deltaRgc{} --- an order of magnitude less, because it declines to act on
four instances in five --- and this is what H4 tests and what it
\textbf{\verdictHFour} on: selective evaluation matches exhaustive evaluation at a
fraction of the compute, mostly because it rarely acts.

The comparison that matters for practice is against the static ensemble, which costs
one evaluation per instance, takes no decision at all, and is the only policy whose
interval excludes zero in the helpful direction on CRPS. Any adaptive correction stage
has to be argued for against that, and on this benchmark the argument cannot be made
on accuracy at any price. The LLM arms are priced separately in
Section~\ref{sec:llmarms}, because their cost is measured in provider calls rather
than in local re-forecasts and the two do not add.


\begin{table}[htbp]
\centering
\caption{Accuracy against computational cost, in action evaluations and LLM calls per
instance.}
\label{tab:cost}
\small
\fittable{\begin{tabular}{lrrr}
\toprule
Policy & MASE & Action evals & LLM calls \\
\midrule
never correct & 1.0090 & 0.00 & 0.00 \\
rgc & 1.0132 & 0.79 & 0.00 \\
random correct & 1.0365 & 0.21 & 0.00 \\
threshold single & 1.0701 & 2.61 & 0.00 \\
always correct & 1.0978 & 9.00 & 0.00 \\
best fixed & 1.0981 & 1.00 & 0.00 \\
\bottomrule
\end{tabular}
}
\end{table}

\begin{figure}[htbp]
\centering
\includegraphics[width=0.62\textwidth]{fig_cost_accuracy}
\caption{Accuracy versus computation. The region of interest is the lower left: equal or
better accuracy for less computation.}
\label{fig:cost}
\end{figure}

\subsection{Where correction helps}

If correction pays anywhere, it should pay where the forecast is under stress.
Figure~\ref{fig:strata} breaks the change down by generator family, which on
synthetic data means a breakdown across known structural conditions rather than across
heterogeneous datasets whose differences are uncontrolled.

The families where correction earns most are the ones where the base forecast is
misspecified in a way a single action repairs; contaminated series, where removing
outliers and refitting recovers the fit, and series whose assumed seasonal period is
wrong. The families where it earns least are those where the base model was already
adequate, and there the intervention is a coin flip with a cost. This is a statement
about the \emph{conditional} value of acting, and it is the motivation for gating at
all; Section~\ref{sec:robustness} tests it directly by perturbing the series along
controlled axes instead of reading it off the families that happen to exist.


\begin{figure}[htbp]
\centering
\includegraphics[width=0.94\textwidth]{fig_stratum_gains}
\caption{Change by generator family. Because synthetic families carry ground-truth
structure, this is a comparison across known conditions rather than across heterogeneous
datasets.}
\label{fig:strata}
\end{figure}

\section{Ablation Studies}
\label{sec:ablations}

Each ablation removes one component and replays against the identical counterfactual
tensor, so every configuration sees the same forecasts and the same action outcomes and
only the decision policy varies. Where a signal family is ablated it is removed from the
estimator's features as well as from the diagnoser, so the gate cannot reintroduce it.

\begin{table}[htbp]
\centering
\caption{Component ablations, replayed against the identical counterfactual tensor.
A positive $\Delta$ means removing the component made results worse. Intervals are
95\% series-level bootstrap against the full system; $\dagger$ marks one excluding
zero. Only one does.}
\label{tab:ablations}
\small
\fittable{\begin{tabular}{lrrrr}
\toprule
Configuration & MASE & $\Delta$ vs. full (\%) [95\% CI] & Corr. rate & Cost \\
\midrule
full & 1.0132 & 0.00~[0.00, 0.00] & 0.20 & 0.79 \\
no self critique & 1.0151 & 0.12~[-2.75, 2.09] & 0.24 & 2.38 \\
no uncertainty & 0.9986 & -1.36~[-3.97, 0.10] & 0.21 & 0.83 \\
no model disagreement & 1.0119 & -0.11~[-0.41, 0.16] & 0.20 & 0.80 \\
no changepoint & 1.0029 & -0.98~[-3.52, 0.42] & 0.20 & 0.82 \\
no anomaly & 1.0108 & -0.18~[-3.01, 1.67] & 0.20 & 0.78 \\
no adaptive tools & 1.0151 & 0.12~[-2.75, 2.09] & 0.24 & 2.38 \\
no backtesting & 1.0404 & 4.41$^\dagger$~[0.92, 7.41] & 0.28 & 0.83 \\
no correction loop & 1.0090 & -0.13~[-3.02, 1.84] & 0.00 & 0.00 \\
\bottomrule
\end{tabular}
}
\end{table}

One component is load-bearing. Removing backtesting costs \ablNoBacktestingSeries\%
([\ablNoBacktestingSeriesCIlo, \ablNoBacktestingSeriesCIhi]\%) and pushes the
intervention rate from \ablFullRate\% to \ablNoBacktestingRate\%: without an
estimate of how the base forecast has actually been doing, the policy loses its grip on
when to leave well alone. That is the same signal the language model was handed and did
not use (Section~\ref{sec:llmarms}).

No other component's removal is distinguishable from zero at the series level. The
uncertainty signals, changepoint proximity, anomaly signals and cross-model disagreement
each have intervals that include zero, and so does the self-critique step. An earlier
draft read the instance-level point estimates in this table as evidence that several
components were mildly harmful; the intervals do not support that, and we no longer say
it. What they do support is weaker: only the crudest signal in the diagnostic apparatus,
recent backtest error, does identifiable work.

The last row is the one to read. Removing the correction loop entirely, so the policy
never acts, changes MASE by \ablNoCorrectionLoopSeries\%
([\ablNoCorrectionLoopSeriesCIlo, \ablNoCorrectionLoopSeriesCIhi]\%). The full
system and no system are indistinguishable, which is the headline result of
Section~\ref{sec:headline} restated as an ablation.

Self-critique and adaptive tool use are the same configuration in this framework, since
the critique decides whether to call a further tool; the two rows are identical for that
reason. Removing it raises compute from \ablFullCompute{} to
\ablNoAdaptiveToolsCompute{} action evaluations per instance while changing accuracy by
\ablNoAdaptiveToolsSeries\% ([\ablNoAdaptiveToolsSeriesCIlo,
\ablNoAdaptiveToolsSeriesCIhi]\%).

\section{Robustness Analysis}
\label{sec:robustness}

Each perturbation axis is applied at increasing severity to the \emph{same} base series,
with forecast targets left untouched, so degradation is attributable to the perturbation
rather than to differences between datasets.

\begin{figure}[htbp]
\centering
\includegraphics[width=\textwidth]{fig_robustness}
\caption{Degradation under controlled perturbation.}
\label{fig:robustness}
\end{figure}

The sweep is the sharpest test of the case for correction, because it manufactures the
conditions correction exists to handle. Across \robComparisons{} comparisons ---
seven perturbation axes, four severities, two acting policies, each against never
correcting on the same series --- \textbf{\robImprovements{} show a significant
improvement} and \robDegradations{} show a significant degradation. A correction stage
that cannot beat inaction under injected anomalies, regime shifts, drift, seasonal
change, added noise, missingness or truncated history has not merely failed to win on
average. It has failed where its own rationale places it.

\begin{table}[htbp]
\centering
\caption{Correction against never correcting under controlled perturbation, at the
strongest severity of each axis. Series-level means with bootstrap intervals over
series; negative favours the policy and $\dagger$ marks an interval excluding zero.
The full sweep over all four severities is in the released results.}
\label{tab:robustness}
\small
\fittable{\begin{tabular}{lllrrrr}
\toprule
Perturbation & Severity & Policy & Series & $\Delta$ vs never & 95\% CI & Win rate \\
\midrule
artificial anomalies & 1.00 & always correct & 84 & 0.55\% & [-3.39, 5.97] & 52.4\% \\
artificial anomalies & 1.00 & rgc & 84 & -1.05\% & [-4.88, 2.33] & 54.2\% \\
gaussian noise & 1.00 & always correct & 84 & -2.86\% & [-6.95, 1.02] & 54.8\% \\
gaussian noise & 1.00 & rgc & 84 & -1.15\% & [-2.77, 0.50] & 57.1\% \\
gradual drift & 1.00 & always correct & 84 & 3.36\%$^\dagger$ & [0.45, 6.68] & 41.1\% \\
gradual drift & 1.00 & rgc & 84 & 0.36\% & [-1.92, 2.60] & 50.0\% \\
regime shift & 1.00 & always correct & 84 & 6.40\% & [-0.56, 13.60] & 46.4\% \\
regime shift & 1.00 & rgc & 84 & -1.18\% & [-4.28, 1.98] & 54.8\% \\
seasonal change & 1.00 & always correct & 84 & -0.49\% & [-5.05, 3.34] & 44.0\% \\
seasonal change & 1.00 & rgc & 84 & 1.39\% & [-0.05, 3.12] & 48.8\% \\
short history & 1.00 & always correct & 84 & 2.43\% & [-3.49, 8.24] & 45.2\% \\
short history & 1.00 & rgc & 84 & -1.57\% & [-4.36, 0.71] & 51.2\% \\
\bottomrule
\end{tabular}
}
\end{table}

\paragraph{H2 holds, in a restricted form worth stating precisely.}
H2 predicted that gains \emph{concentrate} under shift, anomalies and regime change.
Concentration is not level, so the pre-registered test is a difference-in-differences:
a policy's advantage over never correcting in a perturbed stratum, minus its advantage
on the same series unperturbed. Both differences are formed per series before
subtraction, so the pairing survives on each side and the bootstrap resamples one unit.

The difference-in-differences design is not fastidiousness. Perturbation is applied to
the series history as well as to the forecast window, so it moves the MASE denominator
and rescales every error on that series: mean accept-MASE falls from
\robSeverityMild{} to \robSeveritySevere{} severity purely as a consequence. Levels are
therefore not comparable across severities, and only a within-severity contrast against
the same series means anything.

Across \robHTwoTested{} strata, \robHTwoRaw{} concentrate before correction, against
roughly two expected by chance at this many tests, and \robHTwoSurviving{} survive the
Holm correction the analysis plan fixes. H2 is therefore \textbf{\verdictHTwo{} in a
restricted form}: concentration appears for \texttt{\robHTwoPolicies}, under
\robHTwoAxes{}, at the mildest severity applied (\robHTwoSeverity{}), and nowhere else
in the sweep. Two strata out of \robHTwoTested{} is a narrow verdict and we report it as
one.

\paragraph{Neither surviving stratum belongs to the adaptive gate.}
Both are unconditional correction. The conditional value of correcting is real and it
is detectable, and the machinery built to detect it did not detect it. The policy
designed to fire exactly where correction pays did not fire there, while the policy that
fires everywhere captured it. That is a more specific indictment of diagnosis-driven
gating than the aggregate result, and it points the same way as
Section~\ref{sec:decomposition}.

\paragraph{Why the gains vanish as severity rises: we cannot say.}
A natural reading of concentration at mild severity only is that severe damage is
beyond what the menu can repair. We cannot test that here. The obvious measure,
how much an oracle could still recover at each severity, uses an oracle selected on
test outcomes, and Section~\ref{sec:decomposition} shows that such an oracle is almost
entirely winner's curse on this benchmark. The perturbation sweep has no validation
origins, so a realizable version cannot be computed. We therefore leave the severity
pattern unexplained, and note only that it is also what one would expect if the two
surviving strata were partly chance: two of forty-eight, both at one severity level.

\section{Error Analysis}
\label{sec:erroranalysis}

We characterise the calibrated policy's mistakes on CRPS, where it acts on 97\% of
instances. Of these, 47\% improved and 38\% harmed the forecast. Harm is concentrated:
the worst 1\% of corrections account for roughly half of all damage summed across
corrected instances, and a single \texttt{refit\_post\_changepoint} correction on one
instance cost as much as many small wins combined. This heavy tail is why the calibrated gate's design includes an optional risk-averse lower bound on predicted gain. Validation selected that bound's width as zero (deviation~D4), so the deployed gate thresholds the point estimate; the tail is real, but on this benchmark widening the bound did not help on held-out origins.

Two further findings temper any adaptive-policy claim. First, among corrected instances,
the calibrated gain estimate used to \emph{select} the action does not predict which
individual corrections help ($r=-0.02$, $p=0.22$): the calibration is informative about
which actions are worth trying in aggregate, not about which specific instances will
benefit. Second, \texttt{ensemble} accounts for the large majority of the calibrated
policy's corrections and of its benefit, which is consistent with Section~\ref{sec:headline}, where the unconditional ensemble is the only policy that improves CRPS reliably: the adaptive layer spends most of its budget approximating the static policy without matching it.
